\renewcommand{\Anno}{2021}

\renewcommand{\NumeroProblema}{1}

\renewcommand{\Testo}{
Antonio, uscendo da casa per recarsi al lavoro, deve percorrere una strada rettilinea lunga 1.600~m; a metà percorso c'è un semaforo con i soli colori rosso e verde, ognuno dei quali ha una durata di un minuto. Antonio usa come veicolo uno scooter che ha l'accelerazione massima di 1~m/s$^2$ e i freni molto potenti, tali da poter decelerare a piacimento o addirittura arrestare lo scooter istantaneamente se necessario.

\begin{enumerate}
\renewcommand{\labelenumi}{(\alph{enumi})}

\item  Antonio, acceso lo scooter, vede che il rosso è appena comparso.

Determinare con quale condotta di guida arriverà al lavoro il prima possibile.

\item  Recentemente il semaforo è stato dotato di un cronometro che indica tra quanto tempo il colore cambierà.

Determinare, in funzione del colore del semaforo e dell'indicazione data dal cronometro nel momento in cui Antonio accende lo scooter, il minimo tempo che occorre ad Antonio per arrivare al lavoro.

\item  Un giorno Antonio, dopo aver acceso lo scooter, vede che il semaforo è rosso, ma il cronometro è guasto, per cui non è possibile sapere quanto manca al prossimo verde.

Determinare il tempo minimo entro il quale Antonio è \emph{sicuro} di poter arrivare al lavoro, e quale condotta di guida deve tenere per arrivare \emph{con certezza} entro questo tempo massimo.

\end{enumerate}

{\small Nota: supponiamo che lo scooter sia puntiforme, che non ci siano problemi di traffico, che la strada sia a senso unico (per cui non si può tornare indietro), e che né la strada né lo scooter abbiano limiti di velocità. }
}

\renewcommand{\Risposta}{}

\renewcommand{\Aiutino}{}

\renewcommand{\SoluzioneUfficiale}{}

\renewcommand{\SoluzioniAlternative}{}

\renewcommand{\Osservazioni}{}

\renewcommand{\Commenti}{}
